\section{Source audit and proposed integration corrections}
\label{audit:section}

The audit found no new counterexample to a theorem in the inherited
material on which the proofs here depend. It did find several places
where the status of a direction must be updated when the newer reports
are integrated. An older question remains a valid historical record
even when a later incoming article answers it. The appropriate change
is a dated completion note and a link to the proof, while retaining
the original provenance.

\paragraph{A verified editorial duplication.}
The retrieved manuscript \file{README.md} repeats the paragraph beginning
``The complete 1,032-input first-letter X family consecutively.
Remove its second occurrence while retaining the first and the surrounding
research-status text. This duplication was already noted in the latest
incoming source-status report; it is an editorial cleanup, not a new
mathematical erratum. The separate statements about pending canonical
analytic integration remain appropriate and should not be rewritten as
though every incoming result had already been incorporated.

\subsection{What should change in the current research map}

The fixed-gap map should now record complete gaps $2,4,6,8,10$.
The gap-eight proof belongs to \cite{PolygammaResonances}; only gap ten
is completed in this article. The smaller-order-through-twelve results
from earlier work are another source of coverage and should not be
counted again as new gap-ten cases. The pair $(14,24)$ is included as
an explicit example outside that earlier smaller-order range.

The compact inverse-row question receives a constructive bounded-size
answer in Theorem~\ref{gt:inverse}: at fixed half-gap $m$, every entry
is a quotient involving an $(m+1)$-square determinant. This does not
answer the stronger request for a short product or for a recurrence
using only a fixed number of final rows. Nor does the gap-ten proof
remove the finite computational component; the source's request for
a structural proof without finite completion remains open.

The three-cluster question $B^r,C,D$ in
\emph{Gamma Normalizations and Colored Moments} is answered by the
explicit dictionary in Section~\ref{tri:sec:main}
\cite{GammaColored}. A status note should retain the distinction between
the finite colored dictionary for one spectral term and the convergent
infinite spectral sum for the Gamma face. The much broader cluster
pattern $B^r,C^s,D^t$ with arbitrary multiplicities remains open in the
form posed there.

The collision question should now distinguish three levels. The
two-factor theorem for all positive monomial scales and all Stieltjes
indices is inherited. The present arbitrary-factor theorem evaluates
the complete local correction at Stieltjes index zero, uniformly for
all positive paths. Complete scalar values are additionally obtained
for the one-derivative centered family, including the uncentered
triple. The corresponding arbitrary-factor theorem at positive
Stieltjes indices is a further question.

Finally, the first and sixth research questions of the expanded
\emph{Universal Gamma Normal Forms} package receive the answers in
Section~\ref{hp:section}. The universal normal coefficients themselves
remain an attributed input. The new finite formula is for the
compensating polynomial $H_m$, and the new analytic result is the
convergence of the infinite substitutions for a holomorphic field.

\subsection{Conventions that materially affect correctness}

\paragraph{Nonvanishing does not imply positivity.}
The real fundamental determinant has
$f_{6,5}>0$, $f_{b,5}<0$ for $7\le b\le17$, and $f_{18,5}>0$.
Thus the conjectural strengthening $f_{b,5}>0$ at every admissible
index is false. This is an explicit obstruction to that possible
heuristic, not a claim that the predecessor asserted it. The
finite-field certificate proves nonvanishing. Real positivity is
proved separately only for $b\ge73728$.

\paragraph{The constant depends on the scale relation.}
The nonresonant triple in \eqref{intro:triple} cannot be applied to
$(\alpha_1,\alpha_2,\alpha_3)=(3,2,1)$.
The extra terms in \eqref{pc:integer-triple} are a concrete correction
to such an extension. A power--logarithm constant is also different
from the ordinary limit of the original divergent family. The
explicit $Q$ subtraction states exactly which ordinary limit is
available.

\paragraph{Exact regular coefficients are necessary for arbitrary paths.}
Replacing $A_{S,d}(\boldsymbol a)$ by a prematurely truncated expression
in the shifts can destroy uniformity. For example, if
$a_{\min}=e^{-1/t^2}$ and $a_{\max}=t$, then
$a_{\max}|\log a_{\min}|=1/t$ does not tend to zero.
Theorem~\ref{pc:uniform} keeps the shift dependence of every retained
regular coefficient exactly; its remainder proof therefore needs no
assumption on relative rates.

\paragraph{Labelled derivatives and confluent derivatives differ.}
The mixed operator $\partial_{C_1}\cdots\partial_{C_m}$ is applied
before $C_1=\cdots=C_m$ is imposed. A total derivative along that
diagonal has extra multiplicities and other partial derivatives.
The factorial $(m-1)!$ in the confluent divided difference is part
of the theorem. Likewise, the three-cluster primitive must retain
its two edge terms; its mixed derivative alone does not determine
the primitive without boundary data.

\paragraph{Endpoints and branches stay in the formulas.}
The compensator $UH_m$ and the endpoint term $YH_m$ in
\eqref{hp:weighted} are essential. Discarding a total derivative
under an unweighted finite-part integral requires its own endpoint
calculation. For the colored formulas the displayed real branch is
continued using letters outside the integration path. Iterated-integral
shuffle keeps those letters; a stuffle conversion outside the initial
disk needs additional branch bookkeeping if it introduces new
cumulative colors.

\paragraph{The arithmetic questions remain arithmetic questions.}
An invertible rational coefficient matrix does not prove independence
of its evaluated constants. A depth-three representation does not
prove that depth three is necessary. The Gaussian $S_6$ identity and
the current revised $S_8$ period candidate remain unresolved in the
inspected research programme; the present results provide no proof of
those equalities. These restrictions should remain visible in any
combined manuscript overview.

\section{Further research questions}
\label{questions:section}

The following questions are chosen for their potential to produce
identities, structural formulas and exact evaluations. The thresholds
used in the gap-ten proof are tools for nonvanishing, rather than
the main objects of interest.

\begin{enumerate}[label=\textbf{Q\arabic*.},leftmargin=*]
\item \textbf{Arbitrary unequal even harmonic tails.}
Prove or disprove $\det M_{b,m}\ne0$ for every $b>m\ge1$.
Can a transform, a sign-regular kernel, a Pad\'e interpretation or a
spectral argument account for the nonvanishing despite the finite
sign changes? The full two-parameter problem is stronger than
settling each of finitely many gaps.

\item \textbf{The next fixed gap and a reusable exact completion theorem.}
Complete gap twelve using the effective Gamma basis and the causal
convolution method. A useful further result would be an explicit
algorithm that takes $m$, derives a valid analytic threshold and
produces a fully replayable finite certificate, including a principled
choice of primes when one reduction vanishes. A zero residue is not
evidence of a zero rational determinant.

\item \textbf{Closed Green columns and inverse harmonic coordinates.}
Find a hypergeometric, Pad\'e or integrable representation of
$\mathsf G_{s,N}$ that eliminates its long causal recurrence.
Determine whether the bounded determinant in
\eqref{gt:inverse-formula} admits a short quotient of special functions.
The resulting formulas would turn high-order zeta-coordinate sums
into more recognizable identities.

\item \textbf{Positive-index Stieltjes collisions with three or more factors.}
For $q_i\ge0$, extend the uniform local theorem to
$\prod_i\partial_x^{r_i}\gamma_{q_i}(x+a_i)$, with a fully evaluated
correction and a uniform error for arbitrary positive paths.
The singular kernels now contain logarithms of several shifted
variables. An effective route is to introduce independent spectral
powers before differentiating, while retaining every regular
coefficient needed for uniformity. A formal derivative of an
unproved uniform family is insufficient.

\item \textbf{A finite classification of resonant scale contributions.}
For a fixed derivative pattern and fixed number of factors, classify
the integer relations among the $\alpha_i$ that can contribute to
the log-free constant. The rational kernel makes each specified
monomial case finite. Is there a compact chamber description with
explicit transition polynomials and amplitude logarithms?
The difference between \eqref{intro:triple} and
\eqref{pc:integer-triple} provides a first exact test case.

\item \textbf{More than one derivative-bearing factor.}
Find full scalar evaluations for centered products with two or more
nonzero derivative orders. The local correction is already given by
Theorem~\ref{pc:uniform}; the unresolved step is the merged product.
Which patterns are total derivatives of finite polygamma polynomials,
and which force marked-zeta or genuine new moment coordinates?

\item \textbf{Generating functions for the ordered constant rows.}
Use \eqref{pc:all-powers-formula} to sum in $N$ and in the derivative
position $r$. Seek polylogarithmic or Gamma expressions for these
double generating functions, and relations with harmonic-number
weights. The adjacent-difference and zero-sum identities give two
exact constraints on any proposed answer.

\item \textbf{A finite primitive for more displaced masses.}
Integrate \eqref{tri:eq:Gamma-mixed} for
$B^r,C_1,\ldots,C_m$ with $m\ge3$, keeping all lower-dimensional
face data. Is depth at most $m+1$ sufficient uniformly in $r$?
Can the edge reconstruction be expressed by a finite symmetric
dictionary with controlled letters, rather than by unresolved
parameter derivatives?

\item \textbf{Three clusters with arbitrary multiplicities.}
Extend $B^r,C,D$ to $B^r,C^s,D^t$ and identify the smallest useful
uniform depth bound. The one-unit multiplicities in the present
proof produce an elementary simplex resolvent; the general problem
requires retaining additional Dirichlet powers without losing the
finite colored description.

\item \textbf{Arithmetic collapse of the Gamma spectral sum.}
Find special positive $a$ and nonnegative masses for which
\eqref{tri:eq:Gamma} reduces to finitely many Gamma, zeta or
multiple-polylogarithm constants. Coincident masses, rational
shifts and imaginary-Gamma multiplication identities give natural
families to test. Any reduction should state its period basis and
branch conventions before invoking a numerical relation search.

\item \textbf{Signs and full support of the compensators.}
The exact data through weight $24$ suggest that for every $m\ge2$
all monomials of weight $m-1$ occur in $H_m$, with coefficient sign
$(-1)^{m-|\boldsymbol a|}$.
Prove or refute this statement. A positive integral formula for the
absolute coefficients would explain the cancellations in the
Euler-homotopy expression and could also control their denominators.

\item \textbf{Sharp analytic domains and continuation of the canonical series.}
Determine the true convergence domains of $\mathfrak N_U$ and
$\mathfrak H_U$ for $U=\psi$ and $U=-\gamma_q$.
The sufficient domain $|t|<x/2$ makes the translation identity
rigorous; it is not asserted to be sharp. Can the normal and
primitive pieces continue separately to the natural singular locus
of the Gamma or Stieltjes ratio, or do they develop additional
singularities that cancel only in their sum?

\item \textbf{Label-preserving and difference-algebra normal forms.}
Develop a canonical complement for several independent fields that
retains the individual Gamma factors and their mixed derivative
relations. In a separate direction, replace the ordinary derivation
by a difference operator and derive the changed homotopy and product
rules. Simply substituting a $q$-Gamma function into an ordinary
derivative formula would not resolve this structural question.

\item \textbf{The unresolved Gaussian periods.}
Continue the Gaussian $S_6$ and revised $S_8$ problems in a declared
relation system. Seek a new analytic transformation, a complete
iterated-integral reduction or a certified functional identity.
Finite relation searches and numerical agreement can guide that
work but do not become proofs merely by increasing precision or
the size of a relation matrix.
\end{enumerate}
